
\section{FAIR}\fbeg{}\begin{frame}{\ft{\q{Findable, Accessible, Interoperable, Reusable}}}

\begin{ftxt}

According to the popular \q{FAIR} standard, data should be 
Findable, Accessible, Interoperable, and Reusable.  Arguably 
the most important consideration here is accessibility.  
In the typical case, a researcher will learn about 
some scientific work or experiment from peer-reviewed 
literature.  Once they are reading a specific 
publication --- and should they wish to examine 
the research in greater detail --- the question becomes 
how readily they can acquire the associated raw 
data and examine it usefully.\\\vspace{-13pt}

Interoperability means that multiple data sets can be 
studied and analyzed side-by-side, and Reusability 
means that data and code can be reused in new experiments 
--- maybe reproducing and/or extending the original 
data set.  But accessibility is a precondition for everything 
that follows.

\end{ftxt}

\end{frame}

